\documentclass[12pt,letterpaper]{article}

\usepackage[spanish]{babel}
\usepackage[utf8]{inputenc}
\usepackage[T1]{fontenc}
\usepackage{graphicx}
\usepackage{fancyhdr}
\usepackage{geometry}
\usepackage{amsmath, amssymb}
\usepackage{multirow}

\geometry{margin=2.5cm}
\setlength{\headheight}{52pt}


\pagestyle{fancy}
\fancyhf{}

\fancyhead[L]{
    \includegraphics[height=1.5cm]{ur_image.png}
}

\fancyhead[R]{
\begin{tabular}{r}
\textbf{Monitoria Álgebra Lineal} \\
Gabriel Fernando Salcedo Zamora \\
\end{tabular}
}

\renewcommand{\headrulewidth}{0.4pt}

\begin{document}

\section*{Mínimos Cuadrados}

\vspace{0.3cm}

Dado un sistema sobredeterminado $Ax \approx b$, donde no existe solución exacta,
se busca un vector $\hat{x}$ que minimice el error cuadrático:

\[
\min_x \|Ax - b\|^2
\]

La solución de este problema está dada por la ecuación normal:

\[
\hat{x} = (A^T A)^{-1} A^T b
\]

Esta expresión permite encontrar la recta (o modelo lineal) que mejor se ajusta a un conjunto de datos.

El vector solución se interpreta como:
\[
\hat{x} =
\begin{bmatrix}
b \\
m
\end{bmatrix}
\]

donde $m$ es la pendiente y $b$ es la intersección con el eje $y$.  
Así, la recta se construye como:
\[
y = mx + b
\]

\vspace{0.5cm}

\begin{enumerate}
    \item De los siguientes problemas encuentre la recta que mejor se ajusta a los puntos dados:
    \begin{enumerate}
        \item $(5,-3),(2,-4),(-3,6)$
        \item $(1,3),(-2,4),(7,0)$
        \item $(-1, 10), (-2, 6), (-6, 6), (2, -2)$
        \item $(1, 4), (-2, 5), (3, -1), (4, 1)$
    \end{enumerate}
    
    \item Un fabricante compra grandes cantidades de refacciones para cierta máquina. Él encuentra
que este costo depende del número de cajas compradas al mismo tiempo y que el costo
por unidad disminuye conforme el número de cajas aumenta. Supone que el costo es una
función lineal del volumen y de las facturas anteriores obtiene la siguiente tabla:
\begin{table}[h!]
\centering
\begin{tabular}{ |cc| } 
\hline
Número de cajas compradas & Costo total (USD)  \\
\hline
10 & 150 \\
30 & 260 \\
50 & 325 \\
100 & 500 \\
175 & 670 \\
\hline
\end{tabular}
\end{table}
\end{enumerate}

\newpage

\section*{Solución}


\begin{enumerate}
    \item
    \begin{enumerate}
        \item $y = -\dfrac{17}{14}x + \dfrac{9}{7}$
        \item $y = -\dfrac{19}{42}x + \dfrac{68}{21}$
        \item $y = -\dfrac{7}{30}x + \dfrac{79}{15}$
        \item $y \approx -0.88x + 3.57$
    \end{enumerate}

    \item $y \approx 3.05x + 142.6$

\end{enumerate}

\end{document}

